La solución se implementa mediante scripts compatibles con \texttt{/bin/sh}. Todos los recursos de trabajo de la práctica (claves, sales y ficheros originales, cifrados y recuperados) se centralizan en \texttt{scripts/crypto-workspace}. El entorno se prepara con \texttt{scripts/prepare-inputs.sh}; los scripts de cada pregunta se ejecutan desde su directorio. Se empleó OpenSSL~4.1.0, instalado mediante Homebrew en macOS.

\lstinputlisting[firstline=1,lastline=4,caption={Comprobación del entorno de ejecución.},label={lst:p1-entorno-real}]{recursos/salidas/pregunta-1/entorno.txt}

\textbf{Preparación del entorno:} El script central \texttt{scripts/prepare-inputs.sh} crea los directorios comunes, incluidos los destinados a las claves asimétricas y las firmas, y copia el enunciado en \texttt{originals} si no existe. Puede repetirse sin borrar recursos ni generar claves.

\begingroup
\renewcommand{\lstlistingname}{Script}
\setstretch{1}
\lstinputlisting[style=programacompleto,language=bash,
 caption={Preparación del workspace común: prepare-inputs.sh.},
 label={lst:workspace-prepare}]{recursos/listados/prepare-inputs.sh}
\endgroup

\textbf{Generación de claves y sales:}
El script \texttt{gen-keys.sh} trabaja en el workspace preparado y utiliza \texttt{openssl rand} para generar una clave maestra y diez sales de 16 bytes. 
\texttt{umask 077} restringe el acceso a los archivos nuevos al usuario propietario, evitando que otros usuarios del sistema puedan acceder a ellos y modificarlos.
La función \texttt{generate} llama al comprobador antes de confirmar que la inicialización ha finalizado correctamente.
Si ya existen claves o falla su generación, el script solicita escribir \texttt{DELETE} antes de borrar y recrear todo el workspace. Cualquier otra respuesta conserva los archivos y cancela la ejecución.

\Needspace{9\baselineskip}
\begingroup
\renewcommand{\lstlistingname}{Script}
\setstretch{1}
\lstinputlisting[style=programacompleto,language=bash,
 caption={Generación y comprobación de la maestra y las sales: gen-keys.sh.},
 label={lst:p1-gen-completo}]{recursos/listados/pregunta-1/gen-keys.sh}
\endgroup

La ejecución registrada de \texttt{gen-keys.sh} confirmó la validación de la maestra y las diez sales. Las salidas conservan las rutas relativas de la ejecución original; los scripts mostrados utilizan el workspace común:

\begingroup
\renewcommand{\lstlistingname}{Salida y verificación}
\setstretch{1}
\lstinputlisting[style=terminal,basicstyle=\ttfamily\footnotesize,caption={Ejecución de gen-keys.sh con validación satisfactoria.},label={lst:p1-inicializacion-real}]{recursos/salidas/pregunta-1/inicializacion-validada.txt}
\endgroup

A continuación se muestran la clave maestra y las dos primeras sales en hexadecimal, consultadas con \texttt{od} para verificar su contenido visualmente.
Cada par de dígitos representa un byte y cada línea contiene los 16 bytes del archivo.

\begingroup
\setstretch{1}
\lstinputlisting[style=terminal,basicstyle=\ttfamily\footnotesize]{recursos/salidas/pregunta-1/maestra-y-sales-hex.txt}
\endgroup

\textbf{Validación de la generación de claves:} El script \texttt{check-inputs.sh} comprueba que existen tanto la clave maestra como las diez sales, y que cada archivo contiene exactamente 16 bytes.
Después compara los 45 pares de sales mediante \texttt{cmp -s}. Muestra \texttt{PASS} si todas las comprobaciones se superan; cualquier archivo ausente, tamaño incorrecto o sal repetida provoca un error, sin modificar las entradas.

\begingroup
\renewcommand{\lstlistingname}{Script}
\setstretch{1}
\lstinputlisting[style=programacompleto,language=bash,
 caption={Validación de la clave maestra y las sales: check-inputs.sh.},
 label={lst:p1-prepare-completo}]{recursos/listados/pregunta-1/check-inputs.sh}
\endgroup

\textbf{Derivación y separación de claves:} El script \texttt{derive-keys.sh} aplica PBKDF2 con HMAC-SHA-256 y 600\,000 iteraciones para cada sal, manteniendo la misma clave maestra.
Los comandos \texttt{od} y \texttt{tr} representan las entradas binarias en hexadecimal; las opciones \texttt{hexpass} y \texttt{hexsalt} permiten que OpenSSL recupere sus bytes originales \ref{ref:openssl-kdf-cli}.
Cada llamada genera 48 bytes: \texttt{dd} guarda los primeros 16 como clave de cifrado y los 32 restantes como clave de autenticación.

Tras separar cada par, se elimina la clave derivada intermedia y \texttt{wc -c} muestra los tamaños reales de las dos claves generadas para cada fichero.
Las claves se guardan en \texttt{secrets/derived-keys} y \texttt{iterations.txt} conserva el coste empleado (número de iteraciones).
La creación del directorio mediante \texttt{mkdir} impide sobrescribir una derivación existente; \texttt{set -eu} detiene el script ante órdenes fallidas o variables no definidas.

\begingroup
\renewcommand{\lstlistingname}{Script}
\setstretch{1}
\lstinputlisting[style=programacompleto,language=bash,
 caption={Derivación y separación de claves de cifrado y autenticación por cada fichero: derive-keys.sh.},
 label={lst:p1-derive-completo}]{recursos/listados/pregunta-1/derive-keys.sh}
\endgroup

La ejecución de \texttt{derive-keys.sh} produjo la siguiente salida.
Para cada fichero, \texttt{wc -c} confirmó 16 bytes de cifrado y 32 bytes de autenticación, con un subtotal de 48 bytes generados por fichero:

\begingroup
\renewcommand{\lstlistingname}{Salida y verificación}
\setstretch{1}
\lstinputlisting[style=terminal,basicstyle=\ttfamily\footnotesize,caption={Ejecución de derive-keys.sh con validación de tamaños.},label={lst:p1-derivacion-real}]{recursos/salidas/pregunta-1/derivacion.txt}
\endgroup

Como ejemplo, se muestran en hexadecimal las dos claves generadas para el fichero 01.
Cada par de dígitos representa un byte: la clave de cifrado ocupa una línea de 16 bytes y la de autenticación, dos líneas de 16 bytes (32 bytes totales).
La representación permite consultar los valores binarios sin modificar los archivos.

\begingroup
\setstretch{1}
\lstinputlisting[style=terminal,basicstyle=\ttfamily\footnotesize]{recursos/salidas/pregunta-1/claves-derivadas-hex.txt}
\endgroup

Se obtienen diez pares de claves de los tamaños previstos y sales distintas. La maestra, las sales y las iteraciones permiten reproducir la derivación.

\textbf{Generación de los ficheros originales:} El script auxiliar \texttt{gen-files.sh} utiliza \texttt{openssl rand} para generar diez ficheros binarios, sin un contenido específico.
El tamaño inicial es de 1\,024 bytes y aumenta un byte por fichero, hasta 1\,033; así se comprueba el relleno con longitudes distintas. Puede indicarse otro tamaño inicial como argumento.

Los archivos nuevos quedan restringidos al propietario mediante \texttt{umask 077}; si falla alguna generación, se elimina el archivo incompleto de esa operación.

\begingroup
\renewcommand{\lstlistingname}{Script}
\setstretch{1}
\lstinputlisting[style=programacompleto,language=bash,
 caption={Generación de diez ficheros aleatorios con distintos tamaños: gen-files.sh.},
 label={lst:p1-gen-files}]{recursos/listados/pregunta-1/gen-files.sh}
\endgroup

La ejecución creó los diez originales con los tamaños previstos:

\noindent\begin{minipage}{\linewidth}
\begingroup
\renewcommand{\lstlistingname}{Salida y verificación}
\setstretch{1}
\lstinputlisting[style=terminal,basicstyle=\ttfamily\footnotesize,
 caption={Generación de los originales de 1024 a 1033 bytes.},
 label={lst:p1-files-real}]{recursos/salidas/pregunta-1/generacion-ficheros.txt}
\endgroup
\end{minipage}

\textbf{Cifrado y autenticación:} El script \texttt{encrypt-files.sh} asigna a cada fichero las claves derivadas de cifrado y autenticación con el mismo identificador.
Reutiliza las sales y las 600\,000 iteraciones conservadas en la fase de derivación; comprueba la existencia de las entradas y los tamaños de las claves y la sal antes de procesarlas.

Cada salida \texttt{file-XX.protected} comienza con una cabecera $H_i$ de 48 bytes: 4 para la versión \texttt{P1E1} (identifica AES-128-CBC y HMAC-SHA-256), 2 para el identificador, 10 para las iteraciones, 16 para la sal y 16 para un nuevo vector aleatorio.

OpenSSL cifra el original con \texttt{enc -aes-128-cbc}, la clave $K_i^{\mathrm{cif}}$ y el vector generado, añadiendo automáticamente el relleno PKCS\#7. La opción \texttt{-nosalt} evita un encabezado de sal adicional de \texttt{enc}: la sal de PBKDF2 ya figura en $H_i$ y se utilizó al obtener las claves.
Después, \texttt{dgst -sha256 -mac HMAC} calcula la etiqueta con $K_i^{\mathrm{mac}}$ sobre la cabecera y el criptograma, y la añade al final en formato binario \ref{ref:openssl-gcm}.

El fichero protegido contiene, por tanto, $H_i\Vert C_i\Vert T_i$. La salida se escribe directamente en \texttt{encrypted/pregunta-1}, con permisos \texttt{600}; un fallo elimina el contenedor incompleto que se estuviera generando.

\begingroup
\renewcommand{\lstlistingname}{Script}
\setstretch{1}
\lstinputlisting[style=programacompleto,language=bash,
 caption={Cifrado AES-128-CBC y autenticación HMAC: encrypt-files.sh.},
 label={lst:p1-encrypt-files}]{recursos/listados/pregunta-1/encrypt-files.sh}
\endgroup

La ejecución protegió los diez ficheros utilizando sus claves correspondientes:

\noindent\begin{minipage}{\linewidth}
\begingroup
\renewcommand{\lstlistingname}{Salida y verificación}
\setstretch{1}
\lstinputlisting[style=terminal,basicstyle=\ttfamily\footnotesize,
 caption={Cifrado y autenticación de los diez originales.},
 label={lst:p1-cifrado-files-real}]{recursos/salidas/pregunta-1/cifrado-ficheros.txt}
\endgroup
\end{minipage}

Todos los contenedores ocupan 1\,120 bytes. Los originales de 1\,024 a 1\,033 bytes producen un texto con relleno de 1\,040 bytes, al que se suman los 48 de cabecera y los 32 de etiqueta.
En particular, el original de 1\,024 bytes, múltiplo de 16, recibe un bloque completo de relleno. La coincidencia de tamaños no implica que los criptogramas sean iguales.

\textbf{Verificación y descifrado:} El script \texttt{decrypt-files.sh} solo necesita la maestra $K_M$ y los ficheros protegidos. Comprueba el tamaño y la versión del contenedor, el identificador y el formato de las iteraciones.
Antes de derivar, limita estas últimas al intervalo de 1 a 10\,000\,000 para acotar el coste de procesar una cabecera todavía no autenticada.

De cada cabecera recupera $S_i$, $c$ y $V_i$. Con la misma maestra, \texttt{openssl kdf} reconstruye los 48 bytes de PBKDF2-HMAC-SHA-256; \texttt{cut} separa sus primeros 16 bytes y los 32 restantes para recuperar $K_i^{\mathrm{cif}}$ y $K_i^{\mathrm{mac}}$.
La representación hexadecimal contiene dos caracteres por byte. Estas claves se mantienen en memoria, sin leer ni recrear los archivos de claves derivadas \ref{ref:openssl-kdf-cli}.

El script recalcula el HMAC sobre todos los bytes anteriores a la etiqueta almacenada. Recorre sus 32 bytes completos, combina cada pareja mediante XOR y acumula las diferencias mediante OR: solo acepta si el acumulador termina en cero.
Este recorrido evita detenerse en la primera diferencia.

Tras aceptar la autenticación, extrae el criptograma y lo descifra con la clave reconstruida y el vector almacenado. OpenSSL valida y retira el relleno antes de completar la salida en \texttt{decrypted/pregunta-1}.
Un error posterior a su creación elimina el archivo incompleto; un rechazo de autenticación no produce un nuevo fichero descifrado.

\begingroup
\renewcommand{\lstlistingname}{Script}
\setstretch{1}
\lstinputlisting[style=programacompleto,language=bash,
 caption={Reconstrucción PBKDF2, verificación HMAC y descifrado AES-128-CBC: decrypt-files.sh.},
 label={lst:p1-decrypt-files}]{recursos/listados/pregunta-1/decrypt-files.sh}
\endgroup

La ejecución confirmó la reconstrucción con 600\,000 iteraciones, la verificación y la recuperación de los diez ficheros:

\noindent\begin{minipage}{\linewidth}
\begingroup
\renewcommand{\lstlistingname}{Salida y verificación}
\setstretch{1}
\lstinputlisting[style=terminal,basicstyle=\ttfamily\footnotesize,
 caption={Verificación HMAC y descifrado AES-128-CBC de los diez contenedores.},
 label={lst:p1-descifrado-files-real}]{recursos/salidas/pregunta-1/descifrado-ficheros.txt}
\endgroup
\end{minipage}

Finalmente, se comparó cada original con su recuperación mediante \texttt{cmp -s}, ejecutando el bucle desde \texttt{sh}.
Cada mensaje \texttt{PASS} se muestra únicamente si ambos archivos coinciden byte a byte; la primera diferencia detendría las comparaciones:

\noindent\begin{minipage}{\linewidth}
\begingroup
\renewcommand{\lstlistingname}{Salida y verificación}
\setstretch{1}
\lstinputlisting[style=terminal,basicstyle=\ttfamily\footnotesize,
 caption={Comprobación exacta de los diez originales y sus recuperaciones.},
 label={lst:p1-comparacion-files-real}]{recursos/salidas/pregunta-1/comparacion-ficheros.txt}
\endgroup
\end{minipage}

Los diez mensajes \texttt{PASS} confirman la recuperación exacta de todos los originales.
Con ello se completa la generación de diez claves de cifrado diferentes desde la maestra, su aplicación a diez ficheros y la posterior verificación y recuperación de sus contenidos.
